% Page budget: 1.55 pages | Owns: negation, OBC derivation, additional-state scope, implementation, and true-parameter condition.
% WRITING GUIDE
% Purpose: Explain negation, derive the final state-level orthogonal-by-construction map, and delimit its interpretability guarantee.
% Start here: Decisions D-185, D-190, D-192, D-194, D-198, D-202, D-203, D-222, D-229, and D-232 in docs/decisions.md.
% Supporting material: Current state-level OBC derivations under scripts/gantry/orthogonal-by-construction/documentation and their thesis-readiness guidance.
% Mandatory papers: Read Györök's projection-regularization paper and orthogonal-by-construction paper in full before writing; keep their guarantees and assumptions separate.
% Research-plan anchor: Preserve Aspect 3's goal of protecting physical interpretability; state clearly how the final construction differs from the proposed regularisation.
% Current decisions: Reduced identifiable coordinates, data-derived reference tuples, a tangent-only one-step physical-state projection, per-epoch refresh, and no affine comparison arm.
% Choices to justify: Protected parameter coordinates, reference-set construction, projected rows, rank tolerance, exclusion of the affine offset, refresh frequency, and treatment of additional states.
% Implementation audit: Read model_augmentation/fit_systems/obc.py, scripts/gantry/gantry_dynamic/obc_gantry.py, obc_diagnostics.py, and the OBC attachment in model.py before fixing the derivation.
% Reference check: Verify the original projection method and separate its claims from the state-level, nonlinear-parameter, LPV, and additional-state extensions developed here.
% Cautions: docs/orthogonal-projection-plan.md describes an older soft-penalty route; one-step state orthogonality is not trajectory-output orthogonality or guaranteed true recovery.
% Planned figures: New geometric projection figure with a physical and additional-state partition inset; show the remaining uniqueness limitation as clearly as the protected direction.
% Section is finished when: The protected space, coordinate frame, gradient treatment, added-state scope, and true-recovery condition are explicit.
\section{Interpretability by Orthogonal Construction}\label{sec:obc}

Section~\ref{sec:aug_structure} left the split between the baseline and the
network non-unique, which puts the physical meaning of
$\theta_{\mathrm{base}}$ at risk.
This section removes the first-order part of that freedom with an
orthogonal-by-construction (OBC) parametrisation.
Gy\"or\"ok et al.\ introduced OBC for discrete-time input-output models that
are linear in their parameters~\cite{gyorok2026obc}.
They name state-space baselines without full-state measurement as an open
extension~\cite[Sec.~6]{gyorok2026obc}.
The gantry baseline is such a model, and its transition is also nonlinear in
the physical parameters.
Sections~\ref{sec:obc_negation} to~\ref{sec:obc_scope} treat negation, the
protected space, the corrected model, and the scope of its guarantee.

\subsection{Negation in Joint Estimation}\label{sec:obc_negation}
Joint estimation lets the network reproduce a change of the physical
parameters.
In \eqref{eq:aug_transition}, $f_{\mathrm{aug}}$ is added to the physical rows
of the baseline transition.
A change of $\theta_{\mathrm{base}}$ moves these rows along the columns of the
one-step sensitivity~\cite[Eq.~(15)]{gyorok2025l4dc}
\begin{equation}
 \Phi_{\bar\theta}(\tilde x,u)
 =\left.\frac{\partial f_{\mathrm{base}}(\theta_{\mathrm{base}},\tilde x,u)}
   {\partial\theta_{\mathrm{base}}}\right|_{\theta_{\mathrm{base}}=\bar\theta_{\mathrm{base}}}
   \in\R^{6\times10},
\label{eq:obc_sens}
\end{equation}
where $\bar\theta_{\mathrm{base}}$ is the expansion point and
$f_{\mathrm{base}}$ is the normalised RK4 step of
Section~\ref{sec:aug_structure}.
To first order, the network output can cancel a parameter change
$\delta\theta$:
\begin{equation}
\begin{aligned}
 &f_{\mathrm{base}}(\bar\theta_{\mathrm{base}}+\delta\theta,\tilde x,u)
 +\bigl[f_{\mathrm{aug}}(\theta_{\mathrm{aug}},\hat x,u)
 -\Phi_{\bar\theta}(\tilde x,u)\,\delta\theta\bigr]\\
 &\quad\approx f_{\mathrm{base}}(\bar\theta_{\mathrm{base}},\tilde x,u)
 +f_{\mathrm{aug}}(\theta_{\mathrm{aug}},\hat x,u).
\end{aligned}
\label{eq:obc_negation}
\end{equation}
Both sides predict the same next state, so the data cannot attribute
$\delta\theta$ to either component~\cite[Ex.~2]{gyorok2026obc}.
We call this negation.
As a result, the estimate of $\theta_{\mathrm{base}}$ is not unique, and it
can drift to physically unrealistic values~\cite[Sec.~3.1]{gyorok2026obc}.

On the gantry, the negated directions are generalized forces.
Differentiating the solved form of \eqref{eq:eom} with respect to the $j$-th
combination gives
\begin{equation}
 \frac{\partial\ddot q}{\partial\theta_{\mathrm{base},j}}
 =-M(Y)^{-1}\Bigl[\frac{\partial M(Y)}{\partial\theta_{\mathrm{base},j}}\,\ddot q
 +\frac{\partial C}{\partial\theta_{\mathrm{base},j}}\,\dot q
 +\frac{\partial K}{\partial\theta_{\mathrm{base},j}}\,q\Bigr].
\label{eq:obc_gantry_sens}
\end{equation}
The RK4 step carries this acceleration sensitivity into
\eqref{eq:obc_sens}.
The mass combinations act through $\ddot q$, the damping combinations through
$\dot q$, and the stiffness combination through $q$.
A negating network output therefore acts as an added mass, damping or stiffness force.
Through the factor $M(Y)^{-1}$, these directions change with the payload
position.

The construction protects the ten identifiable combinations
$\theta_{\mathrm{base}}$ of Section~\ref{sec:params}, not the fourteen raw
parameters.
Along the four null directions of the raw parameters the dynamics do not
change, so their sensitivity has rank ten.
The coefficient fitted in Section~\ref{sec:obc_map} would then not be unique.
The combinations are trained in the logarithmic and bounded coordinates of
Section~\ref{sec:params}.
The map from these coordinates to $\theta_{\mathrm{base}}$ has a nonsingular
Jacobian, so both give the same range of \eqref{eq:obc_sens}.
The aim is a unique estimate of $\theta_{\mathrm{base}}$.
Whether it equals the true value is a separate condition
(Section~\ref{sec:obc_scope}).

\subsection{Protected Space on Reference Tuples}\label{sec:obc_space}
A penalty can discourage negation but not remove it.
The projection regularisation of~\cite{gyorok2025l4dc} penalises the part of
the network output that lies in the range of the baseline.
Its weight trades model accuracy against orthogonality, and choosing it is
not straightforward~\cite[Sec.~1]{gyorok2026obc}.
OBC removes that part exactly instead.
Consider a baseline that is linear in its parameters $\theta_b$, with
regressor stack $\Phi$ and network output stack
$F^{\mathrm{ANN}}_{\theta_a}$ over the training data.
OBC corrects the network output as~\cite[Eqs.~(8), (10), (11)]{gyorok2026obc}
\begin{equation}
\begin{aligned}
 \theta_{\mathrm{aux}}&=(\Phi^\top\Phi)^{-1}\Phi^\top F^{\mathrm{ANN}}_{\theta_a},\\
 \widetilde F^{\mathrm{ANN}}_{\theta_a}&=F^{\mathrm{ANN}}_{\theta_a}-\Phi\,\theta_{\mathrm{aux}},\\
 0&=\Phi^\top\widetilde F^{\mathrm{ANN}}_{\theta_a},\\
 \hat y_k&=\phi(x_k)\theta_b+f^{\mathrm{ANN}}_{\theta_a}(x_k)-\phi(x_k)\,\theta_{\mathrm{aux}},
\end{aligned}
\label{eq:obc_gyorok}
\end{equation}
where $\theta_{\mathrm{aux}}$ is the least-squares fit of the network output
on the regressor, and $\phi(x_k)$ is the regressor at the input-output
history $x_k$.
The corrected output is orthogonal to the baseline range for every network
weight, so no penalty weight is needed.
After training, the last line predicts at new data with the fitted
coefficient held fixed~\cite[Eq.~(13)]{gyorok2026obc}.

On the gantry, the sensitivity \eqref{eq:obc_sens} takes the place of the
regressor.
The baseline transition has no linear-in-parameters form, because
$\theta_{\mathrm{base}}$ enters through $M(Y)^{-1}$ in every RK4 stage.
Its first-order expansion about $\bar\theta_{\mathrm{base}}$ is linear in
$\theta_{\mathrm{base}}$ with regressor $\Phi_{\bar\theta}$, plus an
offset~\cite[Eq.~(16)]{gyorok2025l4dc}.
The network and the baseline add in the six physical rows of the state
transition, so these rows form the protected stack.
We protect the range of $\Phi_{\bar\theta}$ only, whereas the penalty
of~\cite[Eq.~(18)]{gyorok2025l4dc} also covers the offset.
The offset multiplies no estimated parameter.
By \eqref{eq:obc_negation}, network output along it can therefore not be
traded against a change of $\theta_{\mathrm{base}}$ to first order.

Gy\"or\"ok et al.\ evaluate the regressor at every training sample, which
requires its arguments at every sample.
Here the regressor needs the full physical state, but only the positions are
measured.
States from an open-loop simulation of the
baseline~\cite[Sec.~3]{gyorok2025l4dc} are unsuitable, because such a
simulation drifts on the free axes $X$ and $Y$
(Section~\ref{sec:aug_closed_loop}).
We therefore form a fixed reference tuple $(\tilde x_i,\bar x_i,u_i)$ from
every valid training sample,
\begin{equation}
\begin{aligned}
 q_i&=P^{-\top}y_i,\qquad
 \dot q_i=\frac{-q_{i+2}+8q_{i+1}-8q_{i-1}+q_{i-2}}{12\,T_s},\\
 \tilde x_i&=\col(q_i,\dot q_i),\qquad \bar x_i=0,
\end{aligned}
\label{eq:obc_tuples}
\end{equation}
where $y_i$ is the measured position, $u_i$ the recorded actuator force, and
$i=1,\dots,N$ runs over the samples of all training records.
The fourth-order central difference never crosses a record boundary.
Such an auxiliary evaluation set is permitted
by~\cite[Sec.~3.2]{gyorok2026obc}.

Each tuple is a one-step evaluation point, not a simulated trajectory.
It carries its measured payload position $Y_i$, so the stack contains the
sensitivity at every visited scheduling value.
The recorded input already contains the feedback action, so the controller
does not enter the tuples.
The additional states are not measured, and the tuples set them to zero.
The tuples are normalised in the same way as the training signals and stay
fixed during training.

\subsection{Corrected Augmented Model}\label{sec:obc_map}
On the tuples, the construction makes the physical-row output of the network
orthogonal to the protected range.
Stacking the sensitivity and the network output over the tuples gives
\begin{equation}
\begin{aligned}
 \Phi_{\bar\theta}(\tilde X,U)
 &=\col\bigl(\Phi_{\bar\theta}(\tilde x_1,u_1),\dots,
   \Phi_{\bar\theta}(\tilde x_N,u_N)\bigr)\in\R^{6N\times10},\\
 F_{\mathrm{aug}}
 &=\col\bigl(f_{\mathrm{aug}}(\theta_{\mathrm{aug}},\hat x_1,u_1),\dots,
   f_{\mathrm{aug}}(\theta_{\mathrm{aug}},\hat x_N,u_N)\bigr),\\
 \theta_{\mathrm{aux}}&=\Phi_{\bar\theta}(\tilde X,U)^{+}F_{\mathrm{aug}},\\
 0&=\Phi_{\bar\theta}(\tilde X,U)^\top
   \bigl(F_{\mathrm{aug}}-\Phi_{\bar\theta}(\tilde X,U)\,\theta_{\mathrm{aux}}\bigr),\\
 0&=\Phi_{\bar\theta}(\tilde X,U)^\top\,
   \frac{\partial\bigl(F_{\mathrm{aug}}-\Phi_{\bar\theta}(\tilde X,U)\,
   \theta_{\mathrm{aux}}\bigr)}{\partial\theta_{\mathrm{aug}}},
\end{aligned}
\label{eq:obc_ours}
\end{equation}
where $\hat x_i=\col(\tilde x_i,\bar x_i)$ and ${}^{+}$ denotes the
Moore-Penrose pseudoinverse.
The protected range contains every first-order change of the stacked
physical transition that $\theta_{\mathrm{base}}$ can generate.
The stack must have full column rank, which is checked each time it is
built.
The fourth line is~\cite[Lemma~3]{gyorok2026obc} applied to this stack.
The last line follows from differentiating the fourth, because the stack does
not depend on $\theta_{\mathrm{aug}}$.
As a result, a gradient step on the network cannot move its corrected output
into the protected range.

The corrected model subtracts the fitted parameter change at every rollout
step:
\begin{equation}
\begin{aligned}
 \begin{bmatrix}\tilde x_{k+1}\\ \bar x_{k+1}\end{bmatrix}
 &=\begin{bmatrix}f_{\mathrm{base}}(\theta_{\mathrm{base}},\tilde x_k,\hat u_k)\\ 0\end{bmatrix}\\
 &\quad+\begin{bmatrix}f_{\mathrm{aug}}(\theta_{\mathrm{aug}},\hat x_k,\hat u_k)
   -\Phi_{\bar\theta}(\tilde x_k,\hat u_k)\,\theta_{\mathrm{aux}}\\
   g_{\mathrm{aug}}(\theta_{\mathrm{aug}},\hat x_k,\hat u_k)\end{bmatrix},
\end{aligned}
\label{eq:obc_model}
\end{equation}
where $\hat u_k$ is the closed-loop input of \eqref{eq:aug_controller}.
At each rollout point, the correction evaluates the sensitivity at the
expansion point of the stack and applies the fitted coefficient, as the
last line of \eqref{eq:obc_gyorok} does at new data.
It is the directional derivative of the baseline step in the direction
$\theta_{\mathrm{aux}}$.
The baseline has zero rows for the additional states, so $g_{\mathrm{aug}}$
is not corrected.

The coefficient and the stack are updated at different rates.
The coefficient is recomputed from the current network at every evaluation
of the objective \eqref{eq:aug_loss}, and the gradient passes through this
solve.
Fitting it on the rollout states instead would make the correction at step
$k$ depend on later states of the same window.
The stack and its expansion point are rebuilt only between epochs, at the
current $\theta_{\mathrm{base}}$.
This saves a sensitivity evaluation at every tuple per update, and the range
changes little with the expansion point.
Within an epoch both are fixed, so no gradient reaches
$\theta_{\mathrm{base}}$ through them.
Before every validation and after training, the coefficient is recomputed
from the current network and held fixed for prediction.

\subsection{Scope and Condition for True Parameter Recovery}\label{sec:obc_scope}
The guarantee \eqref{eq:obc_ours} is narrower than orthogonality of the
trained model.
It holds for the one-step transition, stacked over the chosen tuples, and
only locally around $\bar\theta_{\mathrm{base}}$.
It therefore does not establish global uniqueness of $\theta_{\mathrm{base}}$.
It uses the Euclidean inner product $\langle a,b\rangle=a^\top b$ in
normalised state coordinates, so the state normalisation defines its metric.
It is aggregate: only the sum over the tuples vanishes, not each tuple's
term.

The additional states lie outside the guarantee.
The tuples lie on the slice $\bar x=0$, so the coefficient does not see
directions that act only through nonzero additional states.
The update $g_{\mathrm{aug}}$ is never corrected, and it can reach the
physical rows one step later through $f_{\mathrm{aug}}$.
Away from the slice, the rollout still subtracts
$\Phi_{\bar\theta}\theta_{\mathrm{aux}}$, but no stacked orthogonality holds
there.
Consequently, the construction does not make the network's contribution to
the predicted output trajectory orthogonal to the baseline.
The consistency and zero-covariance results
of~\cite[Thms.~16, 18]{gyorok2026obc} therefore do not transfer, because
their proofs use orthogonality of the scored output.

Non-overlap makes the split unique, but it does not make it true.
Let $\theta^*_{\mathrm{base}}$ be the true combinations and
$\hat\theta_{\mathrm{base}}$ their estimate.
The true one-step discrepancy on the tuples is
\begin{equation}
 \Delta^*=\col\bigl(\tilde x_i^{+}-f_{\mathrm{base}}(\theta^*_{\mathrm{base}},
 \tilde x_i,u_i)\bigr)_{i=1}^{N},
\label{eq:obc_delta}
\end{equation}
where $\tilde x_i^{+}$ is the recorded successor of tuple $i$, built with the
same stencil and normalised in the same way.
Suppose that the data are noise-free, that $f_{\mathrm{base}}$ is affine in
$\theta_{\mathrm{base}}$ on a region containing $\hat\theta_{\mathrm{base}}$
and $\theta^*_{\mathrm{base}}$, and that the corrected model reproduces every
recorded successor.
The stack, written $\Phi$, is then the same at every expansion point in
that region.
The stacked one-step identity and the orthogonality of \eqref{eq:obc_ours}
then give
\begin{equation}
 \Delta^*=\Phi\bigl(\hat\theta_{\mathrm{base}}-\theta^*_{\mathrm{base}}\bigr)
 +\widetilde F_{\mathrm{aug}}
 \;\Longrightarrow\;
 \hat\theta_{\mathrm{base}}-\theta^*_{\mathrm{base}}=\Phi^{+}\Delta^*,
\label{eq:obc_bias}
\end{equation}
where $\widetilde F_{\mathrm{aug}}=F_{\mathrm{aug}}-\Phi\theta_{\mathrm{aux}}$
is the corrected network output.
The estimate is therefore exact if and only if $\Phi^\top\Delta^*=0$, which
is the state-level counterpart of~\cite[Cond.~4, Thm.~7]{gyorok2026obc}.
The training objective is a closed-loop multistep output error, so
\eqref{eq:obc_bias} is a one-step surrogate, not the bias of the trained
model.

On the gantry, the recovery condition can hold only through cancellation
over the tuples.
A missing generalized force acts through the acceleration, like the
directions of \eqref{eq:obc_gantry_sens}.
Where the gantry moves, this acceleration sensitivity generically has rank
three, so a nonzero missing force is not orthogonal to it at a single tuple.
Its contributions must therefore cancel over the stack.
In general, this requires a dedicated experiment design, which is possible
when the structure of the missing dynamics is
known~\cite[Sec.~4.1, Remark~5]{gyorok2026obc}.

Two overlap measures, defined with the same inner product, test the
condition and the overlap that remains unprotected.
For a stacked vector $v$ over the tuples, they are
\begin{equation}
\begin{aligned}
 \cos\angle\bigl(v,\ran\Phi\bigr)
 &=\frac{\norm{\Phi\Phi^{+}v}}{\norm{v}},\\
 \cos\angle\bigl(v,F_{\mathrm{base}}\bigr)
 &=\frac{|\langle F_{\mathrm{base}},v\rangle|}
   {\norm{F_{\mathrm{base}}}\,\norm{v}},
\end{aligned}
\label{eq:obc_overlap}
\end{equation}
where $\Phi$ is the stack at $\bar\theta_{\mathrm{base}}$ and
$F_{\mathrm{base}}=\col\bigl(f_{\mathrm{base}}(\bar\theta_{\mathrm{base}},
\tilde x_i,u_i)\bigr)_{i=1}^{N}$ is the stacked nominal baseline transition.
Applied to $\Delta^*$ with the stack at $\theta^*_{\mathrm{base}}$, the first
measure is zero exactly when the recovery condition holds.
It needs $\theta^*_{\mathrm{base}}$, so it is available in simulation only.
Applied to the corrected network output on held-out tuples, it shows how far
the guarantee extends beyond the reference tuples.
The second measure gives the overlap with $F_{\mathrm{base}}$, which the
construction leaves unprotected.
Section~\ref{sec:results} tests whether the split is start-independent and,
where the condition holds, whether the true combinations are recovered.
